\chapter{おわりに}
\label{ch:conclusion}

\section{本研究のまとめ}

本研究は、極大クリーク列挙問題における Eppstein アルゴリズムの並列化を、単一の実装で終わらせず、測定の連鎖として進めてきた。
ここでその連鎖を振り返る。

出発点は逐次アルゴリズム同士の比較である。
十種類のグラフのうち、実データで最も重い soc-sign-epinions では、Tomita らのピボット版 CLIQUES が167.3秒であるのに対し、縮退順序を用いる Eppstein は94.9秒だった。
疎なグラフほど Eppstein が有利という傾向は、縮退数 $d$ とグラフサイズ $n$ の比 $d/n$ でおおむね説明できた。
この結果を受けて、以降の並列化は Eppstein を対象とし、最も重い soc-sign-epinions を主力のベンチマークに定めた。

次に試したのが、最外層の頂点を worker に配るだけの素朴な並列化である。
バッチサイズ512では、worker 数を4以上に増やしても67〜70秒から縮まなかった。
バッチサイズを64に変更しただけで、同じ8 worker が27.7秒まで縮む。
この一つの変更で結果が大きく動いた事実は、並列化の遅さの原因が worker 数ではなくタスクの配り方にあることを示していた。

そこで負荷分布を頂点単位で測定すると、soc-sign-epinions では上位1\%の頂点が全体の94.7\%の時間を占めていた。
偏りの発生源は縮退順序の末尾（疎な外周から剥がしていくため、密なコアが必然的に末尾に残る）である。
一方 ca-HepPh は最も重い1頂点だけで13\%を占め、どう配ってもその1タスクの分割なしにはスピードアップの上限が7.7倍程度に制限される。
どちらも「並列化が効かない」と要約できてしまいそうだが、前者はタスクの配り方で救える偏りであり、後者は部分問題内部を分割しない限り解決しない別種の問題である。

実行時間をロード、縮退順序計算、プロセスプール起動、並列本体の4区間に分解すると、プール起動の固定費が worker 数にほぼ比例して伸びることが分かった（soc で w=1 の0.268秒から w=8 の1.273秒）。
あわせて worker ごとの稼働率を見ると、w=8 では idle が30.4\%に達し、最も忙しい worker と最も暇な worker の差は36.2秒対11.8秒だった。
この idle は、負荷分布の測定で見つけた偏りが worker の遊休として現れたものである。

配分戦略そのものを block、reversed、interleave の三通りで比較すると、interleave は idle を32.6\%からほぼ0.1\%まで解消した。
ところが compute の短縮は4.3\%にとどまり、idle がほぼ消えたにもかかわらず理論上期待される約1.5倍の短縮には遠く及ばなかった。
idle の解消と引き換えに、全 worker の稼働時間の合計は139秒から198秒へ膨張しており、遊休を埋める効果を稼働時間そのものの膨張が打ち消していた。

この膨張の機構を、taskpolicy による P/E コア固定実験で切り分けた。
同じ計算を P コアで実行すると96.0秒、E コアに固定すると363.6秒であり、この条件では E コアは P コアの4分の1程度の速度しか出ない（固定に用いた background QoS はクロックも制限するため、この比は下限値である）。
メモリをほとんど使わない純計算ループでは速度比がさらに開く（5.75倍）ことから、この差の主因はクロック制限であり、クリーク列挙で差が縮むのはメモリ待ちの時間が低クロックの影響を受けないためだと判断できる。
Oakley による powermetrics 測定 \cite{oakley2024} を参照してクロック比から E コアの実力を推定すると、w=8 の実測スループット（P コア5.26個分相当）は、この推定モデル（$4 + 4\times0.43 \approx 5.7$ P コア分、効率92\%）とおおむね整合した。
配分戦略比較で観測した busy 時間の膨張は、この非対称コアが主因だと測定によって裏づけられたことになる。

最後に、並列化のもう一つの固定費であるメモリコピーの削減に取り組んだ。
multiprocessing の起動方式を spawn から fork に変えると、copy-on-write によってグラフの pickle 転送が不要になり、startup は全グラフで0.002〜0.031秒まで縮んだ。
これにより、spawn では並列化がすべて損だった小さいグラフのうち4つが、fork では w=4 で約2倍の得に転じた。
グラフを CSR 形式の配列として shared\_memory に置く shm 方式は、startup の削減幅では fork に劣るものの、compute 自体が fork より一貫して速く、soc-sign-epinions の w=8 では total が26.5秒に達した。
逐次版の94.9秒（別セッションの測定のため参考値）から数えると約3.6倍にあたり、この時点が本研究の到達点である。

\section{得られた知見}

一連の測定から言えるのは、並列化の性能はアルゴリズムの選択だけで決まるのではないということである。
本研究の範囲では、少なくとも四つの要因が性能に効いていた。
プロセスプール起動とデータ転送にかかる\textbf{起動固定費}、縮退順序の末尾に偏る\textbf{負荷分布}、P コアと E コアで速度が大きく異なる（下限実測で約4倍、全力推定でも2倍以上）\textbf{ハードウェアの非対称性}、そして pickle 転送か copy-on-write か共有メモリかという\textbf{メモリ共有方式}である。

これらは独立に効くのではなく、互いを覆い隠す形で作用していた。
配分戦略の改善で負荷分布起因の idle をほぼ消しても、その分の稼働時間が非対称コアの存在によって膨張し、期待した高速化は得られなかった。
逆に言えば、負荷分布だけを見て「配分は改善済みだからこれ以上は望めない」と判断していたら、非対称コアという別の要因を見落としたままになっていた。
本研究で複数の切り分け実験を積み重ねたのは、この種の要因の重なりを一つずつ分離するためである。

ただし、この知見は Apple M4 搭載の1台のマシン、Python の multiprocessing、十種類のグラフという条件のもとで得られたものであり、他のハードウェアや実装、より大規模なグラフでも同じ配分で効くとは限らない。
特に非対称コアの影響は、P コアと E コアを持たない一般的なサーバ環境では前提から成り立たない。
ここで得られた一般化は、あくまで本研究の実測が支える範囲にとどまる。

\section{今後の課題}

本研究の測定の過程で、実施に至らなかった検証や、切り分けが済んでいない機構がいくつか残っている。

\begin{itemize}
  \item \textbf{copy-on-write 説の切り分け}：fork の compute が shm より遅い観測について、CoW の遅延コスト（参照カウント書き換えによるページコピー）が原因だとする説を、\texttt{gc.freeze()} を使った再測定で検証すること。
  \item \textbf{E コアの実クロックの確定}：E コアの速度低下がクロック制限によるという推定を、\texttt{sudo} 権限下の \texttt{powermetrics} による実クロックの直接観測で確定させること。
  \item \textbf{共有資源競合の内訳の切り分け}：P コアのみの並列でも効率が79\%に落ちる競合について、キャッシュ、メモリ帯域、クロック引き下げのどれによるものかを切り分けること。
  \item \textbf{熱ドリフトの機構の確認}：長時間の連続測定で同じ純計算ループが32\%遅くなる現象の原因（クロック低下と見られるが未確認）を、実クロックと温度の記録で確認すること。
  \item \textbf{work stealing とオーバーデコンポジション}：探索木を worker 数よりずっと多いサブツリーに分割し、暇な worker が動的にキューから取る方式を導入し、静的な配分戦略の限界を超えられるかを検証すること。
  \item \textbf{グラフの CSR 化}：shm 方式で導入した CSR 形式による表現を、他の起動方式や、より大きなグラフに対しても一貫して適用すること。
  \item \textbf{ca-HepPh 型の単一重部分問題の内部分割}：縮退順序の1頂点だけが全体時間の大半を占めるグラフに対して、その部分問題自体をさらに分割する方法を検討すること。
  \item \textbf{fork と共有配列の両立}：起動固定費を fork 並みに抑えつつ、shm 並みの compute 速度も得られる構成を検討すること。
\end{itemize}
